% \documentclass[sigconf, anonymous=true]{acmart}
% \usepackage{graphicx}


% % \usepackage{fullpage}

% % \usepackage{parskip}
% \usepackage{enumitem}
% \usepackage{changepage}
% \usepackage[T1]{fontenc}
% \usepackage{hyperref}
% % \usepackage[margin=1in]{geometry}
% \usepackage{graphicx} 
% \usepackage{xcolor}
% % \usepackage{amsmath}
% % \usepackage{amsthm}
% % \usepackage{amssymb}
% % \usepackage{newtxtext,newtxmath}
% % \usepackage{algorithm}
% \usepackage[ruled]{algorithm2e}
% \DontPrintSemicolon
% % \usepackage{algpseudocode} 

% % \usepackage{algorithm2e}

% % \usepackage{setspace} 
% \usepackage{tikz}
% \usetikzlibrary{mindmap, trees}
% % \onehalfspacing 
% % \algrenewcommand\algorithmiccomment[1]{\hfill\textit{#1}}


% \SetKwInput{KwParam}{Public Parameters}
% \newcommand{\revone}[1]{\textcolor{blue}{#1}}
% \newcommand{\revtwo}[1]{\textcolor{red}{#1}}
% \newcommand{\revthr}[1]{\textcolor{orange}{#1}}
% \theoremstyle{definition} 
% \theoremstyle{definition}  % 设置定义的样式
% \newtheorem{theorem}{Theorem}[section] 
% \newtheorem{definition}{Definition}[section]
% \newtheorem{proposition}{Proposition}[section]
% \newtheorem{lemma}{Lemma}[section]
% \newtheorem{corollary}{Corollary}[section] 
% \newtheorem*{observation}{Observation}
% \newtheorem{example}{Example}[section]
% \newtheorem*{remark}{Remark}

% \newcommand{\bbd}[1]{\textbf{\textcolor{blue}{#1}}}
% \newcommand{\obd}[1]{\textbf{\textcolor{orange}{[TODO: #1]}}}
% \newcommand{\rbd}[1]{\textbf{\textcolor{red}{[QUESTION: #1]}}}
% \newcommand{\com}[1]{{\small \textcolor{gray}{[COMMENT: #1]}}}

% \renewcommand{\today}{\number\day\ \ifcase\month\or January\or February\or March\or April\or May\or June\or July\or August\or September\or October\or November\or December\fi\ \number\year}

% \newcommand{\err}{\mathrm{Error}} 
% \title{Domain Partitioning for Shuffle-DP}
% \author{Anonymous Author}

% \begin{abstract}
%     % Despite the great success of shuffled differential privacy (shuffle DP) in improving the performance of distributed private data release over traditional local DP (LDP), the intractable nature of the shuffler between the data owner and the analyzer introduces additional challenges in terms of both communication cost and utility degradation. 
%     % Thus, the shuffled DP model remains underexplored in various settings, including: 
%     % a) the user setting, where each party may own an unknown number of records; 
%     % b) the dynamic setting, where data arrives as streams and private data is released continuously; and 
%     % c) the personalized setting, where each party may have different privacy needs.
%     % These yet to answer open problems significantly narrow the practical applicability of the shuffled DP model.
%     % \textcolor{blue}{Recent research in shuffled DP has introduced the idea of ``domain partitioning" to achieve almost instance-optimal utility for certain fundamental queries by privately determining key instance-specific parameters of the queries. 
%     % However, this idea has broader applicability beyond just partitioning data values. 
%     % In this work, we show that the concept of ``domain partitioning" can also serve as the foundation of general frameworks to address all of the aforementioned settings, or even their combinations.
%     % We demonstrate that our framework is sufficiently general to enable any existing standard DP mechanism to support both union-preserving and non-union-preserving queries under different partitioning strategies. 
%     % We analyze the framework theoretically and show that it achieves a polylogarithmic utility overhead in the one-round setting, where information is sent only from the data owners, and a doubly logarithmic utility overhead in the multi-round setting, where information can also be sent back, compared to the standard DP mechanism in the central model.}
% \end{abstract}

% \begin{document}

% \maketitle



% % \begin{itemize}
% %     \item User Shuffle-DP: 1. Union--single round okay 2. non-union--distinct count---single round, then log (1-1, 1-2, 1-4...), multi-round, then no need to allocate log privacy budget (1-1, 2-2, 3-4, 5-8), just have two rounds would be fine
% %     \item User Shuffle Dynamic DP--run log instances 1-1, 1-2, 1-4... log binary trees, use the last one that is over a threshold (continual SP)
% %     \item User Shuffle PDP = partition ep/k, then invoke the previous method, single round, then have to run log instances over 1-1, 1-2, 1-4, 1-8
% %     \item Any
% % \end{itemize}



% {
%     % \input{sec/intro}
%     \input{sec/Preliminary}
%     % \subsection{Challenges and Naive Solutions}


% \label{sec:Challenges and Naive Solutions}
% The core difficulty of user-level shuffle-DP arises from variation in user contributions.
% Fundamentally, shuffle-DP relies on $n$ users collaboratively constructing a single global noise distribution. Standard shuffle-DP mechanisms typically assume uniform contributions to calibrate this collective noise. However, the variation of user contributions disrupts this assumption, leading to two challenges: on one hand, an insufficient noise scale fails to protect ``heavy'' contributors, compromising the user-level DP guarantee; on the other hand, an excessive scale based on the worst-case scenario severely degrades utility, making instance-specific error bounds unattainable.
% Although for simple additive queries such as $Q_{\text{sum}}$, we can effectively reduce the problem to the record-level setting by letting users locally aggregate their records into a single record, this idea fails to generalize to queries such as $Q_{\text{range}}$, $Q_{\text{median}}$, or $Q_{\text{hist}}$. 
% These queries require specific distributional information that is irreversibly lost when compressing a user's dataset into a single record.
% % please check, maybe you can take an example to show that this cannot used in range counting...
% % For simple additive queries such as $Q_\text{count}$ and $Q_\text{sum}$, this variability can be easily addressed: users can locally aggregate their records into a single record, thus reducing the problem to the record-level setting.
% % please check: Below we introduce several naive solutoins to with such heterogeneous contribution sizes
% Therefore, to support general queries under user-level shuffle-DP, the protocol must accommodate multiple records per user. 
% % \obd{Can elaborate a bit more upon the problem that is specific to shuffle DP: better use the structure ``to support general ... the protocol must bring the gap between the challenge introduced by multiple records owned by a user and ... existing shuffle DP frameworks ..."}
% However, solving this problem faces fundamental challenges regarding both privacy and utility, as we show below.

% \paragraph{Challenge 1: Privacy}
% \obd{Add one more sentence to summarise what this challenge is and replace ``Privacy with it"}

% At first glance, a natural idea is to let users calibrate their noise adaptively to accommodate varying contribution sizes $m_i$.
% Specifically, a user $i$ contributing $m_i$ randomizes each record with privacy budget $\varepsilon_i = \varepsilon / m_i$, thus achieving overall $\varepsilon$-DP by group privacy (Lemma~\ref{lem:Group Privacy}).
% However, this strategy fundamentally invalidates the privacy guarantees of the shuffle model. 
% As illustrated in Example~\ref{exm:Bit counting problem}, shuffle-DP relies on the principle that all $n$ users collectively amortize a global noise. When users scale their noise individually according to their $m_i$, the total noise accumulated after shuffling may be insufficient to protect the user with the strictest privacy requirement (i.e., the largest $m_i$). 
% Take $Q_\text{count}$ as an example. Consider a scenario where only a single user $j$ contributes $m$ records, while each of the remaining $n - 1$ users contributes only one record. 
% Under the rule $\varepsilon_i = \varepsilon / m_i$, these $n-1$ users would inject noise scaled by sensitivity of $1$, resulting in only $O(1/\varepsilon)$ noise in total.
% This is far too small to mask the contribution of user $j$, which requires a total noise of $O(m/\varepsilon)$.
% Therefore, to ensure user-level shuffle-DP, the protocol must employ a unified noise scale sufficient for the worst-case $m_{\max}(D)$.
 
% % please check: where is the best to introduce that use eps/M can reduce user-DP to record-DP.
% \paragraph{Challenge 2: Global bound-based approaches suffer from poor utility}
% \obd{Add one more sentence to follow the flaw of the previous section.}
% Given the requirement for a unified noise scale for the worst-case $m_i$, a straightforward solution is to adopt the global upper bound $M$ and require all users to pad their datasets with dummy records to reach size $M$. Under this strategy, each user randomizes every record with privacy budget \(\varepsilon' = \varepsilon / M\), thereby guaranteeing uniform protection across users regardless of their actual \(m_i\).
% However, this solution provides prohibitively poor utility, since the error scales with the theoretical worst-case bound $M$ rather than the instance-specific maximum contribution in the dataset, i.e., $m_{\max}(D)$.
% Revisiting the $Q_\text{count}$ example, in an extremely sparse scenario where each user contributes only a single record (i.e., $m_i = 1$ for all $i$), the optimal noise scale is $O(1/\varepsilon)$. However, the global padding strategy still incurs noise of $O(M/\varepsilon)$, resulting in an unnecessary and substantial degradation in accuracy. 
% Consequently, while global padding satisfies the privacy constraint, it imposes an unacceptable utility.
% This motivates the need for a more refined approach that adaptively identifies an appropriate threshold $m_\tau$ to bound user contributions, thereby achieving near-optimal utility while preserving DP.

%     \input{sec/Methodology}
%     \input{sec/Dynamic_user_shuffle_DP}
%     \input{sec/Application}
%     \input{sec/Optimization}
%     \input{sec/Experiment}
%     % \section{Preliminaries}
Let $\mathcal{X}$ be the universe of all possible records.  
A dataset $D=\{x_{1},x_{2},\ldots,x_{n}\}\in\mathcal{X}^{n}$ comprises $n\in\mathbb{N}$ such records.


\subsection{Hypergrid Domain Partitioning}
In this section, we formally define the most general domain partitioning scheme—namely, a hypergrid over the data universe.
We also provide intuition on its properties and advantages.

\begin{definition}[Hypergrid Domain Partitioning (HDP)]\label{def:HDP}
Let $k\in\mathbb{N}$ denote the number of dimensions, and let $\mathbf{L}=(L_{1},\dots,L_{k})$ and $\mathbf{U}=(U_{1},\dots,U_{k})$ be vectors of lower and upper bounds
with $0<L_{j}<U_{j}$ for all $j$.  
For each $j$, set
\begin{equation*}
    N_{j}= \left\lceil \log_{2}\!\left( \frac{U_{j}}{L_{j}} \right) \right\rceil,
    \qquad
    \Lambda_{j}= \{0,\dots,N_{j}-1\},
    \qquad
    \Lambda=\Lambda_{1}\times\cdots\times\Lambda_{k}.
\end{equation*}
For each index vector $\boldsymbol{\ell}=(\ell_{1},\dots,\ell_{k})\in\Lambda$, define
\begin{equation*}
    \mathcal{X}_{\boldsymbol{\ell}}
    =\Bigl\{x\in\mathcal{X}:
        2^{\ell_{j}}L_{j}\le v_{j}(x)<2^{\ell_{j}+1}L_{j}
        \text{ for all }j\Bigr\},
\end{equation*}
where $v_{j}(x)$ denotes the value of $x$ in the $j$-th dimension.
\end{definition}


By construction, the family $\{\mathcal{X}_{\boldsymbol{\ell}}\}_{\boldsymbol{\ell}\in\Lambda}$ forms a hypergrid partition of $\mathcal{X}$ with cardinality $|\Lambda|=\prod_{j=1}^{k}N_{j}=O(\log^{k})$.
This enables any existing DP mechanism to be hypergrid-parallelized as follows.

\begin{definition}[Mechanism Hypergrid Parallelization (MHP)]
Let $\{\mathcal{X}_{\boldsymbol{\ell}}\}_{\boldsymbol{\ell}\in\Lambda}$ be the partition from Definition~\ref{def:HDP}.  
The hypergrid-parallelization operator $\operatorname{MHP}_{\{\mathcal{X}_{\boldsymbol{\ell}}\}_{\boldsymbol{\ell}\in\Lambda}}:(\mathcal{X}^{n}\!\to\!\mathcal{Y})\longrightarrow(\mathcal{X}^{n}\!\to\!\mathcal{Y}^{\Lambda})$
acts on any mechanism $\mathcal{M}:\mathcal{X}^{n}\!\to\!\mathcal{Y}$, where $\mathcal{Y}$ is the output space of $\mathcal{M}$, by
\begin{equation*}
    \bigl[
        \operatorname{MHP}_{\{\mathcal{X}_{\boldsymbol{\ell}}\}_{\boldsymbol{\ell}\in\Lambda}}
        (\mathcal{M})
    \bigr](D)
    :=\bigl(
            \mathcal{M}\bigl(D\cap\mathcal{X}_{\boldsymbol{\ell}}\bigr)
      \bigr)_{\boldsymbol{\ell}\in\Lambda}.
\end{equation*}
\end{definition}

Intuitively, the output of a hypergrid-parallelized mechanism is an
$N_{1}\times\cdots\times N_{k}$ tensor.  
Because the grids are pairwise disjoint, the parallel-composition theorem guarantees that if $\mathcal{M}$ is $(\varepsilon,\delta)$-DP, then $\operatorname{MHP}_{\{\mathcal{X}_{\boldsymbol{\ell}}\}_{\boldsymbol{\ell}\in\Lambda}}(\mathcal{M})$ is also $(\varepsilon,\delta)$-DP.
%     % \section{User DP in the Shuffle Model}
Suppose there are $n$ users and the $i$-th user has $m_i$ records, $\{x_{i1}\ldots, x_{im_i}\}$, where $k_i\leq K$ for all users and $(\Sigma_{j=1}^{m_i} x_{ij}) \in [U]$. 
We assume no prior knowledge about the value of $M$ and $U$, thus they can be set to be conservatively large.

\section{Union Preserving}
When our target is a union-perserving query, meaning that 
The target is to estimate the total sum of the records value of all users
\begin{equation*}
    \text{Sum}(D)=\sum_{i=1}^{n}\sum_{j=1}^{m_i}x_{ij}
\end{equation*}
under the shuffle model.

We partition the domain of the number of records to 
\begin{equation*}
    [1,1], [2,2], [3,4],\ldots,[M/2, M],
\end{equation*}
then for each domain, we run the sum algorithm in [CCS:SumDP] as follows


\begin{algorithm}[t]
\LinesNumbered
\caption{SumUserShuffleDP (Randomizer)}
\label{alg:user_shuffle_sum_randomizer}
\KwIn{$\{x_{i1}\ldots, x_{im_i}\}$ from the $i$-th user, $\varepsilon$, $\delta$, $\beta$, $n$, $U$, $M$}
\KwOut{$\{\{\widetilde{Q}_{\text{sum},\, ijk}\}_{j=1}^{\log U}\}_{k=1}^{\log M}$}

\For{$k \gets \{0,\, 1, \, 2, \, \ldots, \, \log K\}$}{

    $S_i=\Sigma_{j=1}^{m_i} x_{ij}$
    
    $\{\widetilde{Q}_{\text{sum},\, ijk}\}_{j=1}^{\log U} \gets \text{Randomrizer of SumDP}\left(S_i\cdot \mathbf{I}\left(m_i \in [2^{k-1}+1, 2^{k)}]\right)), \frac{\varepsilon}{2^k},\delta,n,U\right)$
}

\textbf{Send:} $\{\widetilde{Q}_{\text{count},\, i}(r_j)\}_{i=1}^{\lceil \log(\hat{\varepsilon}/\check{\varepsilon})\rceil}$
\end{algorithm}


\begin{algorithm}[t]
\LinesNumbered
\caption{SumUserShuffleDP (Analyzer)}
\label{alg:user_shuffle_sum_analyzer}
\KwIn{$\Bigl\{ R^{[2^{j-1}+1,2^{j}]} = \bigcup_i S_i^{[2^{j-1}+1,2^{j}]} \Bigr\}_{j\in\{0,1,2,\ldots,\log U\}},\;
       \varepsilon,\; \delta,\; \beta,\; n,\; U$}
\KwOut{$\widetilde{\mathrm{Sum}}(D)$}

$\tau \gets 0$\;

\For{$j \gets \{0,\,1,\,2,\,\ldots,\,\log U\}$}{
    $\widetilde{\mathrm{Sum}}\!\bigl(D \cap [2^{j-1}+1,2^{j}]\bigr)
        \gets \displaystyle\sum_{y\in R^{[2^{j-1}+1,2^{j}]}} y$\;

    \tcp*[l]{Set $\tau = 2^{j}$ for the last $j$ satisfying (4)}
    \If{$\widetilde{\mathrm{Sum}}\!\bigl(D \cap [2^{j-1}+1,2^{j}]\bigr)
        > 1.3 \cdot 2^{j} \cdot \ln\!\bigl( 2(\log U + 1) / \beta \bigr) / \varepsilon$}{
        $\tau \gets 2^{j}$\;
    }
}

$\widetilde{\mathrm{Sum}}(D)
   \gets \displaystyle\sum_{j\in\{0,1,2,\ldots,\log \tau\}}
          \widetilde{\mathrm{Sum}}\!\bigl(D \cap [2^{j-1}+1,2^{j}]\bigr)$\;

\Return{$\widetilde{\mathrm{Sum}}(D)$}
\end{algorithm}



%     % \input{sec/straw}
%     % \input{sec/count}
%     % \input{sec/framework}
%     % \input{sec/local}
%     % \input{sec/shuffle}
%     % \input{sec/sum}
%     % \input{sec/max}
    
% }

% \newpage

% \bibliographystyle{acm}
% \bibliography{ref}

% \clearpage
\appendix
\renewcommand{\thetheorem}{4.\arabic{theorem}}
\setcounter{theorem}{0} 
\section{Proof of Theorem~\ref{the: static two-round}}

\label{sec:Proof of static two}


\begin{theorem} 
    Given any $\varepsilon > 0, \delta >0, n \in \mathbb{Z}_+, $ $U \in \mathbb{Z}_{+}$, $M \in \mathbb{Z}_{+}$, for any query $Q$ and any base shuffle-DP protocol $\mathcal{P}_Q$, our two-round protocol satisfies the following properties in the static setting:

    \begin{itemize}
    \item The messages received by the analyzer preserve $(\varepsilon,\delta)$-DP;
    \item The total error of this protocol is bounded by:
    \[
    \begin{aligned}
    &\mathrm{\gamma}_{\ell_p}\!\Big(
        Q,\;
        \frac{16}{\varepsilon}\ln\frac{2(\log M +1)}{\beta} \cdot m_{\max}(D)
    \Big)
    \\
    &+\;
    \err_{\ell_p}\!\Big(
        \mathcal{P}_Q,\;
        \frac{\varepsilon}{4m_{\max}(D)},\;
        \frac{\delta}{4m_{\max}(D)},\;
        \frac{\beta}{2},\;
        n\cdot 2m_{\max}(D)
    \Big)
    \end{aligned}
    \]



    \item In expectation, each user sends $1+ o(1)+\mathrm{Msg}(\mathcal{P}_Q, \allowbreak \frac{\varepsilon}{4m_{\max}(D)}, \allowbreak \frac{\delta}{4m_{\max}(D)}, \allowbreak n\cdot 2m_{\max}(D))$ messages.

\end{itemize}

\end{theorem}

\begin{proof}
    For privacy, the privacy budget $(\varepsilon, \delta)$ is split evenly between the $m_\tau$ estimation task and the query evaluation task. 
For the first round, invoking Lemma~\ref{lem: Bit Counting Protocol}, we have that for every $j \in \{0,1,2,\dots,\log M\}$, subdomain $[2^j+1,2^j]$ preserves $(\varepsilon/2,\delta/2)$-DP. Given that each $m_i$ only has an impact on a single subdomain $[2^j+1,2^j]$, 
parallel composition (Lemma~\ref{lem:Parallel Composition}) applies across the disjoint subdomains $\{[2^j+1,2^j]\}_{j\in\{0,1,2,\dots,\log M\}}$, ensuring this round satisfies $(\varepsilon/2,\delta/2)$-DP.
% it follows that the collection $\{[2^j+1,2^j]\}_{j\in\{0,1,2,\dots,\log M\}}$ preserves $(\varepsilon/2,\delta/2)$-DP by parallel composition (Lemma~\ref{lem:Parallel Composition}). 
For the second round, the protocol $\mathcal{P}_Q$ preserves $(\varepsilon/2,\delta/2)$-DP. Finally, by sequential composition (Lemma~\ref{lem:Sequential Composition}), our two-round protocol preserves $(\varepsilon,\delta)$-DP.
% please check that use lemma x.x or xx theory (lemma x.x)

For utility, let $Q(D)$ be the true query result on original dataset $D$, $Q(D^\tau)$ be the true query result on the standardized dataset $D^\tau$, and $\tilde{Q}(D^\tau)$ be our final noisy output. By the triangle inequality, the overall error $|\tilde{Q}(D^\tau) - Q(D)|$ can be decomposed into two parts:
$$|\tilde{Q}(D^\tau) - Q(D)| \le \underbrace{|Q(D^\tau) - Q(D)|}_{\text{Bias}} + \underbrace{|\tilde{Q}(D^\tau) - Q(D^\tau)|}_{\text{Noise}}.$$
\textit{Bounding the Bias}: 
% The bias is caused by users whose contribution $m_i$ exceeds $m_\tau$. 
The bias stems from the truncation of real records during the contribution standardization process. Let $m'$ denote the total number of records discarded from users with $m_i > m_\tau$. Since the bias $|Q(D^\tau) - Q(D)|$ is bounded by $\gamma_{\ell_p}( Q, m')$, our primary objective is to derive an upper bound for $m'$.
Let $C_j$ and $\tilde{C}_j$ denote the true and estimated number of users in the subdomain $I_j = [2^{j-1}+1, 2^j]$, respectively. Let $m_\tau = 2^l$. The algorithm selects $l$ as the largest index passing the threshold $T= \frac{2}{\varepsilon}\ln\frac{2(\log M +1)}{\beta}$. This implies that for all $j > l$, the estimated count $\tilde{C}_j$ did {not} pass the threshold, i.e., $\tilde{C}_j \le T$. 
Lemma~\ref{lem: Bit Counting Protocol} guarantees that for any such subdomain $I_j$ ($j > l$), with probability at least $1-\frac{\beta}{2(\log M +1)}$, the estimation error is bounded by $|\tilde{C}_j - C_j| \le T$. Combining these inequalities ($C_j \le \tilde{C}_j + |\tilde{C}_j - C_j|$), for any $j > l$, the true count is bounded by:
% Lemma~\ref{lem: Bit Counting Protocol} ensures that for any unselected subdomain $I_j$ ($j > l$), with probability at least $1-\frac{\beta}{2(\log M +1)}$, we have 
% \[
% |\tilde{C}_j-C_j| \le \frac{2}{\varepsilon}\ln\Big(\frac{2(\log M +1)}{\beta}\Big).
% \]
% Consequently, for any $j > l$, the true count is bounded by:
$$ C_j \le \frac{4}{\varepsilon}\ln\Big(\frac{2(\log M+1)}{\beta}\Big).
% = O\Big(\frac{1}{\varepsilon}\ln\frac{\log M}{\beta}\Big).
$$
Now, we calculate the total number of excluded records $m'$. A user in subdomain $I_j$ (where $j > l$) contributes at most $2^j$ records but is clipped to $m_\tau = 2^l$. The loss is at most $2^j - 2^l$.
Note that for $j > K = \lceil \log m_{\max}(D) \rceil$, the true count $C_j = 0$. Thus, summing over $j$ from $l+1$ to $K$:
$$\begin{aligned}
m'
&= \sum_{j=l+1}^{K} C_j \cdot (2^j - 2^l) \\
&\le \frac{4}{\varepsilon}\ln\frac{2(\log M +1)}{\beta} \cdot \sum_{j=l+1}^{K} (2^j - 2^l) \quad  \\
% &\le \frac{4}{\varepsilon}\ln(\frac{\log M+2}{\beta}) \cdot \sum_{j=l+1}^{K} 2^j \\
% &= \frac{4}{\varepsilon}\ln(\frac{\log M+2}{\beta}) \cdot (2^{K+1} - 2^{l+1})\\
% &\le \frac{4}{\varepsilon}\ln\frac{2(\log M +1)}{\beta} \cdot 4 m_{\max}(D)\\
% \\
% & = O\left(m_{\max}(D)\cdot\frac{1}{\varepsilon}\log\frac{\log M}{\beta}\right) 
&\le \frac{4}{\varepsilon}\ln\frac{2(\log M +1)}{\beta} \cdot 2^{K+1} \\
&\le \frac{16}{\varepsilon}\ln\frac{2(\log M +1)}{\beta} \cdot m_{\max}(D).
\end{aligned}$$
This explicitly validates the property (a) claimed in section~\ref{sec: static round 1}, confirming that the number of excluded records is bounded by $\tilde{O}(m_{\max}(D)/\varepsilon)$.
Consequently, the bias is bounded by: $$\mathrm{\gamma}_{\ell_p}\big(Q, \frac{16}{\varepsilon}\ln\frac{2(\log M +1)}{\beta} \cdot m_{\max}(D)\big).$$
with probability $1-\frac{\beta}{2}$ by union bound.



\textit{Bounding the Noise:} 
The noise magnitude depends on the threshold $m_\tau$. Specifically, the base protocol $\mathcal{P}_Q$ is invoked with privacy budget $(\frac{\varepsilon}{2m_\tau},\frac{\delta}{2m_\tau})$ and dataset size $n \cdot m_\tau$. To bound the noise, it suffices to prove that with high probability, $m_\tau \le 2m_{\max}(D)$.
Still using $K = \lceil \log m_{\max}(D) \rceil$, for any subdomain index $j > K$, the subdomain $[2^{j-1}+1, 2^j]$ strictly exceeds $m_{\max}(D)$, implying the true count $C_j=0$. By Lemma~\ref{lem: Bit Counting Protocol}, the estimated count $\tilde{C}_j$ is pure noise. Therefore, with probability at least $1-\frac{\beta}{2(\log M +1)}$, $\tilde{C}_j$ does not exceed the error bound $T$.
% please check if it exactly 1-beta/2
By union bound, with probability at least $1-\beta/2$, the algorithm will not select any $j > K$. Therefore, the selected threshold satisfies $m_\tau \le 2^K \le 2 m_{\max}(D)$.
{This confirms property (b).}
% By lemma~\ref{lem: Bit Counting Protocol}, if a subdomain $[2^{j-1}+1,2^j]$ is empty, i.e., $C_j=0$, with probability $1-\frac{\beta}{2(\log M +1)} $, the estimated count $\tilde{C}_j \le \frac{2}{\varepsilon}\ln \frac{2(\log M +1)}{\beta}$. Thus, by union bound, with probability at least $1-\beta/2$, for all $j> m_\tau$, $[2^{j-1}+1,2^j] $ is empty. Therefore, we can guarantee that with probability at least $1-\beta/2$, $m_\tau\le 2\cdot m_{\max}(D)$. As a result, the noise caused by $\mathcal{P}_Q$ is bounded by $\err_{\ell_p}(\mathcal{P}_Q,\frac{\varepsilon}{4m_{\max}(D)},\frac{\delta}{4m_{\max}(D)},\frac{\beta}{2},   n\cdot 2m_{\max}(D))$ with probability at least $1-\beta/2$. 
Substituting this bound into the error function, the noise is bounded by $\err_{\ell_p}(\mathcal{P}_Q,\frac{\varepsilon}{4m_{\max}(D)},\frac{\delta}{4m_{\max}(D)},\frac{\beta}{2}, n\cdot 2m_{\max}(D))$ with probability at least $1-\frac{\beta}{2}$.

Combining bias and noise yields the final error bound:
    \[
    \begin{aligned}
    &\mathrm{\gamma}_{\ell_p}\!\Big(
        Q,\;
        \frac{16}{\varepsilon}\ln\frac{2(\log M +1)}{\beta} \cdot m_{\max}(D)
    \Big)
    \\
    &+\;
    \err_{\ell_p}\!\Big(
        \mathcal{P}_Q,\;
        \frac{\varepsilon}{4m_{\max}(D)},\;
        \frac{\delta}{4m_{\max}(D)},\;
        \frac{\beta}{2},\;
        n\cdot 2m_{\max}(D)
    \Big)
    \end{aligned}
    \]
    with probability at least $1-\beta$.

    For the communication cost, recall that our protocol consists of two rounds, each consuming privacy parameters $(\varepsilon/2,\delta/2,\beta/2)$. 
    In the first round, each user executes the bit-counting protocol \(\mathcal{P}_{Q_\text{count}}\)~\cite{ghazi2021differentially} on each subdomain $I_j$ and parallel composition (Lemma~\ref{lem:Parallel Composition}) applies across the \(\log M+1\) disjoint subdomains.
    Crucially, since these subdomains are disjoint, each user contributes a value of $1$ to exactly one subdomain, which we denoted as $I^*$,  and $0$ to all remaining $\log M$ subdomains.
    According to Lemma~\ref{lem: Bit Counting Protocol}, the expected message cost for $I^*$ is $1 + O(\frac{\log(1/\delta)}{\varepsilon n})$, while when input value is $0$, the protocol only sends out $o(\frac{\log(1/\delta)}{\varepsilon n})$ message.
    Consequently, the total expected communication per user in the first round is:
    \[
        1 \cdot \Big(1 + O\big(\frac{\log(1/\delta)}{\varepsilon n}\big)\Big) + \log M \cdot O\big(\frac{\log(1/\delta)}{\varepsilon n}\big) = 1 + o(1),
    \]
    assuming $n \gg \log M \cdot \log(1/\delta)/\varepsilon$.
    % please check if this inequation is right
    
    % By lemma~\ref{lem: Bit Counting Protocol}, for each subdomain, the expected number of messages sent by a user is bounded by $1+O(\frac{\log (1/\delta)}{\varepsilon n}) = O(1)$. Since there are \(O(\log M)\) subdomains, the total communication per user in this round is $O(\log M)$. 
    In the second round, the base protocol \(\mathcal{P}_Q\) is executed on the standardized dataset with privacy budget $(\frac{\varepsilon}{2m_\tau},\frac{\delta}{2m_\tau})$. Thus, the communication cost is then $\mathrm{Msg}(\mathcal{P}_Q,\frac{\varepsilon}{2m_\tau},\frac{\delta}{2m_\tau},n\cdot m_\tau)$. From the utility analysis, we have \(m_\tau \le 2 m_{\max}(D)\), 
    % please check, with probability xx?
    and therefore this cost is bounded by $\mathrm{Msg}(\mathcal{P}_Q,\frac{\varepsilon}{4 m_{\max}(D)},\frac{\delta}{4 m_{\max}(D)},n \cdot 2m_{\max}(D))$. Combining these two rounds, the total number of messages per user is bounded by:
    $$1+o(1)+ \mathrm{Msg}(\mathcal{P}_Q, \allowbreak \frac{\varepsilon}{4m_{\max}(D)}, \allowbreak \frac{\delta}{4m_{\max}(D)}, n\cdot 2m_{\max}(D)).$$


\end{proof}

\section{Proof of Theorem~\ref{the: static one-round}}

\setcounter{theorem}{1}
\begin{theorem}
    Given any $\varepsilon > 0, \delta >0, n \in \mathbb{Z}_+, $ $U \in \mathbb{Z}_{+}$, $M \in \mathbb{Z}_{+}$, under the static setting, our one-round protocol achieves that:

    \begin{itemize}
    \item The messages received by the analyzer preserve $(\varepsilon,\delta)$-DP;
    \item The total error of this protocol is bounded by:
    \[
    \begin{aligned}
    &\mathrm{\gamma}_{\ell_p}\!\Big(
        Q,\;
        \frac{16}{\varepsilon}\ln\frac{2(\log M+1)}{\beta} \cdot m_{\max}(D)
    \Big)
    +\;
    O\Big(\err_{\ell_p}\!(
        \mathcal{P}_Q,\;    \\
    &
        \frac{\varepsilon}{m_{\max}(D)\cdot \log M},\;
        \frac{\delta}{m_{\max}(D)\cdot \log M},\;
        \frac{\beta}{\log M},\;
        n\cdot m_{\max}(D)
    )\Big)
    \end{aligned}
    \]



    \item In expectation, each user sends $1+o(1)+\sum_{j \in \{0,1,\dots, \log M\}} \allowbreak \mathrm{Msg}(\mathcal{P}_Q, \frac{\varepsilon}{2\cdot 2^j(\log M+1)}, \frac{\delta}{2\cdot 2^j(\log M+1)},n \cdot 2^j)$ messages.

\end{itemize}

\end{theorem}

\begin{proof}
     For privacy, the privacy budget $(\varepsilon, \delta)$ is split evenly between the $m_\tau$ estimation task and the query evaluation task.
    For the $m_\tau$ estimation,since the subdomains $\{I_j\}$ are disjoint, parallel composition (Lemma~\ref{lem:Parallel Composition}) guarantees that this step satisfies $(\varepsilon/2, \delta/2)$-DP.
    For query evaluation, by applying sequential composition (Lemma~\ref{lem:Sequential Composition}) over the $\lceil \log M \rceil + 1$ candidate query results, each with budget $(\frac{\varepsilon}{2(\log M+1)}, \frac{\delta}{2(\log M+1)})$, the total cost sums to $(\varepsilon/2, \delta/2)$-DP.
    By basic composition of these two parts, the overall protocol satisfies $(\varepsilon, \delta)$-DP.

    For utility, similar to the proof of Theorem~\ref{the: static two-round}, the total error is bounded by the sum of the clipping bias and the DP noise.
    Observe that the procedure for determining $m_\tau$ (including the privacy budget allocation and the threshold condition) is algorithmically equivalent to the first round of the two-round protocol. Consequently, the analysis for the number of excluded records holds directly. With probability at least $1-\beta/2$, the clipping bias is bounded by:
    \[
    \mathrm{\gamma}_{\ell_p}\!\Big(
        Q,\;
        \frac{4}{\varepsilon}\ln(\frac{2(\log M+1)}{\beta}) \cdot 4 m_{\max}(D)
    \Big)
    \]
    Regarding the selected threshold, the guarantee from Theorem~\ref{the: static two-round} remains valid: with probability at least $1-\beta/2$, we have $m_\tau \le 2m_{\max}(D)$.
    However, unlike the two-round protocol, the privacy budget for query evaluation is split among $\log M + 1$ candidate subdomains. For the selected query result corresponding to $m_\tau$, the base protocol $\mathcal{P}_Q$ is applied with parameters:
    \[
    \varepsilon' = \frac{\varepsilon}{2(\log M+1) m_\tau}, \quad \delta' = \frac{\delta}{2(\log M+1) m_\tau}, \quad n' = n \cdot m_\tau.
    \]
    Substituting the upper bound $m_\tau \le 2m_{\max}(D)$ into the error expression $\mathrm{Error}_{\ell_p}(\mathcal{P}_Q, \varepsilon', \delta', \beta', n')$, the noise is bounded by:
    \[
    \begin{aligned}
        O\Big( \mathrm{Error}_{\ell_p}\!\big(
        &\mathcal{P}_Q,\;
        \frac{\varepsilon} {m_{\max}(D) \cdot \log M},\;
        \frac{\delta}{m_{\max}(D) \cdot\log M},\; \\
        &\frac{\beta}{\log M},\;
        n \cdot m_{\max}(D)
    \big) \Big).
    \end{aligned}
    \]
    Combining the bias and noise bounds yields the final overall error bound.

For the communication cost, our one-round protocol executes two tasks in parallel: $m_\tau$ estimation and query evaluation across all candidate thresholds.

For the $m_\tau$ estimation task, the analysis is identical to the first round of the two-round protocol. Each user executes the bit-counting protocol $\mathcal{P}_{Q_\text{count}}$~\cite{ghazi2021differentially} on each subdomain $I_j$ with budget $(\varepsilon/2, \delta/2)$. Since the subdomains are disjoint, each user contributes a value of $1$ to exactly one subdomain and $0$ to all others. By Lemma~\ref{lem: Bit Counting Protocol}, the expected message cost is:
\[
1 \cdot \Big(1 + O\big(\frac{\log(1/\delta)}{\varepsilon n}\big)\Big) + \log M \cdot O\big(\frac{\log(1/\delta)}{\varepsilon n}\big) = 1 + o(1),
\]
assuming $n \gg \log M \cdot \log(1/\delta)/\varepsilon$.

For the query evaluation task, each user must send messages for all $\log M + 1$ candidate thresholds $\{2^j\}_{j=0}^{\log M}$. For each subdomain $j$, the user constructs a standardized dataset $D’$ of size $2^j$ and applies the base randomizer $\mathcal{R}_Q$ with privacy budget $(\frac{\varepsilon}{2 \cdot 2^j (\log M+1)}, \frac{\delta}{2 \cdot 2^j (\log M+1)})$ and participation size $n \cdot 2^j$. The communication cost for subdomain $j$ is therefore $\mathrm{Msg}(\mathcal{P}_Q, \frac{\varepsilon}{2\cdot 2^j(\log M+1)}, \frac{\delta}{2\cdot 2^j(\log M+1)}, n \cdot 2^j)$.

Summing over all $\log M + 1$ subdomains, the total expected communication per user is:
\[
1 + o(1) + \sum_{j=0}^{\log M} \mathrm{Msg}\Big(\mathcal{P}_Q, \frac{\varepsilon}{2\cdot 2^j(\log M+1)}, \frac{\delta}{2\cdot 2^j(\log M+1)}, n \cdot 2^j\Big).
\]

\end{proof}

   
\section{Proof of Theorem~\ref{the: dynamic setting guarantee}}

\renewcommand{\thetheorem}{5.\arabic{theorem}}
\setcounter{theorem}{0} 
\begin{theorem}
    Given any $\varepsilon > 0, \delta > 0$, any $T \in \mathbb{Z}_+$, and $n \in \mathbb{Z}_+$, $M \in \mathbb{Z}_+$, for any query $Q$, at all timestamps $t \in \{1, 2, \dots, T\}$ simultaneously, our protocol for the dynamic setting achieves the following guarantees:
\begin{itemize}
    \item The messages received by the analyzer preserve $(\varepsilon,\delta)$-DP;
    \item With probability at least $1-\beta$, the total error is bounded by:
    \[
    \begin{aligned}
    &O\Big(\gamma_{\ell_p}\big( Q, \frac{\log ^{1.5} T \cdot \log M}{\varepsilon} \ln \frac{T \cdot \log M}{\beta} \cdot m_{\max}(D^t)\big)\Big) + O\Big(\err_{\ell_p}\big(\mathcal{P}_Q,
    \\
    & \frac{\varepsilon}{m_{\max}(D^t) \cdot \log M}, \frac{\delta}{m_{\max}(D^t) \cdot \log M}, \frac{\beta}{\log M}, n \cdot m_{\max}(D^t)\big)\Big)
    \end{aligned}
    \]
    \item In expectation, each user sends $(\log M+1)(1+o(1)) + \sum_{j=0}^{\log M} \sum_{k=1}^{2^j} \mathrm{Msg}(\mathcal{P}_Q, \frac{\varepsilon}{2\cdot 2^j(\log M+1)}, \frac{\delta}{2\cdot 2^j(\log M+1)}, n \cdot 2^j)$ messages over the entire time horizon $T$.
\end{itemize}
\end{theorem}


\begin{proof}
    For privacy, the privacy budget $(\varepsilon, \delta)$ is split evenly between the two parallel tasks: dynamic counting and query evaluation, each receiving $(\varepsilon/2, \delta/2)$.

For \emph{dynamic counting}, the user participates in $\log M+1$ interval counting mechanisms. Since the intervals $\{I_j\}$ are disjoint (a user belongs to exactly one interval at any time), we apply parallel composition (Lemma~\ref{lem:Parallel Composition}) over the intervals, allocating budget $(\varepsilon', \delta') = (\frac{\varepsilon}{2(\log M+1)}, \frac{\delta}{2(\log M+1)})$ per interval.
Within each interval $j$, the binary mechanism operates over $\log T+1$ tree levels. By applying sequential composition (Lemma~\ref{lem:Sequential Composition}) over these levels, each level receives budget $(\frac{\varepsilon'}{\log T+1}, \frac{\delta'}{\log T+1})$.
Thus, dynamic counting satisfies $(\varepsilon/2, \delta/2)$-DP.

For \emph{query evaluation}, each user processes their data under all $\log M+1$ candidate thresholds. These thresholds are not disjoint (threshold $j+1$ covers threshold $j$), requiring sequential composition over the $\log M+1$ intervals with budget $(\varepsilon', \delta')$ per interval.
For each interval $j$, the user contributes $2^j$ streams, each running the base protocol with budget $(\frac{\varepsilon'}{2^j}, \frac{\delta'}{2^j})$. By composition over the $2^j$ streams, interval $j$ satisfies $(\varepsilon', \delta')$-DP.
Under the continual observation model~\cite{dwork2010continual,chan2011private}, incremental updates to streams incur no additional privacy cost, as the guarantee applies to the user's complete sequence $D_i = \{x_{i,t}\}_{t=1}^T$.
Thus, query evaluation satisfies $(\varepsilon/2, \delta/2)$-DP.

By basic composition of the two tasks, the overall protocol satisfies $(\varepsilon, \delta)$-DP.


For utility, similar to the static one-round protocol, the total error is bounded by the sum of clipping bias and DP noise.



\emph{Clipping Bias.}
Similar to the proof of Theorem~\ref{the: static two-round}, by Lemma ..., at each time $t$, the true count of interval $I_j$ is bounded by:
\[
C_j \le O(\frac{\log^{1.5} T}{\varepsilon'}\ln \frac{T}{\beta'}).
\]
Let $l = j^*$ and $K = \lceil \log m_{\max}(D^t) \rceil$. We can get the total number of discarded records is:
\[
\begin{aligned}
m'
&= \sum_{j=l+1}^{K} C_j \cdot (2^j - 2^l) \\
&\le O\Big(\frac{\log ^{1.5} T}{\varepsilon'} \ln \frac{T}{\beta'}\Big) \cdot \sum_{j=l+1}^{K} (2^j - 2^l) \\
&\le O\Big(\frac{\log ^{1.5} T \cdot \log M}{\varepsilon} \ln \frac{T\log M}{\beta}\Big) \cdot m_{\max}(D^t).
\end{aligned}
\]




\emph{DP Noise.}
For the selected threshold $m_\tau^t = 2^{j^*}$, the base protocol $\mathcal{P}_Q$ is applied with parameters:
\[
\varepsilon' = \frac{\varepsilon'}{2^{j^*}} = \frac{\varepsilon}{2(\log M+1) \cdot 2^{j^*}}, \quad \delta' = \frac{\delta}{2(\log M+1) \cdot 2^{j^*}}, \quad n' = n \cdot 2^{j^*}.
\]
Substituting the upper bound $2^{j^*} \le 2m_{\max}(D^t)$, the noise is bounded by:
\[
\begin{aligned}
    O\Big( \err_{\ell_p}\big( &\mathcal{P}_Q, \frac{\varepsilon}{m_{\max}(D^t) \cdot \log M}, \frac{\delta}{m_{\max}(D^t) \cdot \log M}, \\ &\frac{\beta}{\log M}, n \cdot m_{\max}(D^t)\big) \Big).
\end{aligned}
\]

Combining the bias and noise bounds yields the stated error guarantee.

For communication cost, the protocol executes two tasks in parallel: dynamic counting and query evaluation.

For \emph{dynamic counting}, each user runs $\log M+1$ independent binary mechanism instances (one per interval). 
Since intervals are disjoint, each user contributes to exactly one interval at any given time. 
By Lemma~\ref{lem: Bit Counting Protocol}, each binary mechanism instance with parameters $(\frac{\varepsilon'}{\log T+1}, \frac{\delta'}{\log T+1})$ sends $1 + o(1)$ messages in expectation.
Over $\log M+1$ intervals, the total cost is:
\[
(\log M+1) \cdot (1 + o(1)) = (\log M+1)(1+o(1)).
\]

For \emph{query evaluation}, the user maintains $\sum_{j=0}^{\log M} 2^j = 2^{\log M+1}-1$ streams across all intervals.
Each stream $S[j][k]$ runs an instance of the base protocol $\mathcal{R}_Q$ with budget $(\frac{\varepsilon'}{2^j}, \frac{\delta'}{2^j})$ and participant count $n \cdot 2^j$.
The communication cost for stream $S[j][k]$ is $\mathrm{Msg}(\mathcal{P}_Q, \frac{\varepsilon'}{2^j}, \frac{\delta'}{2^j}, n \cdot 2^j)$.
Summing over all streams:
\[
\sum_{j=0}^{\log M} \sum_{k=1}^{2^j} \mathrm{Msg}\Big(\mathcal{P}_Q, \frac{\varepsilon}{2\cdot 2^j(\log M+1)}, \frac{\delta}{2\cdot 2^j(\log M+1)}, n \cdot 2^j\Big).
\]

Combining both tasks, the total expected communication per user is:
\[
\begin{aligned}
    &(\log M+1)(1+o(1)) \\&+ \sum_{j=0}^{\log M} \sum_{k=1}^{2^j} \mathrm{Msg}\Big(\mathcal{P}_Q, \frac{\varepsilon}{2\cdot 2^j(\log M+1)}, \frac{\delta}{2\cdot 2^j(\log M+1)}, n \cdot 2^j\Big).
\end{aligned}
\]

\end{proof}


% \end{document}
